\documentclass[11pt]{article}

\usepackage[margin=0.5in]{geometry}
\usepackage{xcolor}
\usepackage{enumitem}
\usepackage{cite}
\renewcommand{\citepunct}{,}
\usepackage{hyperref}
\usepackage{amsmath,amssymb}
\usepackage{amsthm}

\newcommand{\blue}[1]{\textcolor{blue}{#1}}
\newcommand{\red}[1]{\textcolor{red}{#1}}
\newcommand{\rigorousstationarity}{at least a rigorous stationarity relationship}


\newtheorem{proposition}{Proposition}
\newtheorem{example}{Example}

\hypersetup{
  colorlinks=true,
  linkcolor=blue,
  citecolor=blue,
  urlcolor=blue
}

\title{Appeal for Reconsideration of Manuscript}
\author{}
\date{}
\begin{document}
\maketitle

\vspace{0.5em}
\noindent\textbf{Article title:} ``DualHash: A Stochastic Primal-Dual Algorithm with Theoretical Guarantee for Deep Hashing''

\vspace{0.5em}
\noindent\textbf{Manuscript ID:} M187198

\vspace{0.5em}
\noindent Dear Editor-in-Chief,

We respectfully ask for reconsideration of the decision on our manuscript
\#M187198, ``DualHash: A Stochastic Primal-Dual Algorithm with Theoretical
Guarantee for Deep Hashing.'' We sincerely appreciate the time and effort the reviewers and the Associate Editor devoted to evaluating our work. However, we find that some of the concerns underlying the rejection decision may have arisen from misunderstandings of the manuscript’s formulation and theoretical results, partly because certain aspects were not presented with sufficient clarity in the original submission. These concerns do not affect the core contributions of our work and can be fully addressed through appropriate clarifications and revisions. Below, we respond to the reviewers’ major concerns and outline the corresponding revisions we propose to make. 


\section{Reviewer 1: Comments on the Stationarity Relationship}

\vspace{0.5em}
\noindent\textbf{Main comment.} \emph{%
The key question is whether this transformed problem provides meaningful guarantees for the original nonconvex W-regularized hashing objective. I do not find this connection convincingly established. In particular, the paper should clarify whether there is strong duality, zero duality gap, or at least a rigorous stationarity relationship between the transformed primal--dual problem and the original objective.}

\vspace{0.5em}
\noindent\textbullet\quad\textbf{{Our response.} }
We thank the reviewer for this constructive comment. In the submitted manuscript, Lemma~5.2 already establishes a stationarity condition linking the transformed primal--dual problem and the original objective. 
Namely, when a zero-gap primal--dual optimal pair $(x^*,B^*,\Lambda^*)$
exists, $(x^*,B^*,\Lambda^*)\in\operatorname{crit}\mathcal{L}$. 

Before proceeding, we recall the
notation used in the manuscript. Problem~(P) reads
\[
\tag{P}
\min_{x,B}\ J(x,B)=F(x,B)+h(B),
\]
where $h(B)=\lambda\sum_{i,j}\bigl||B_{ij}|-1\bigr|$ is the nonconvex and
nonsmooth W-regularizer.
The dual problem~(D) is obtained by introducing the splitting $Y=B$
and eliminating $Y$ via the Fenchel conjugate~$h^*$:
\[
\max_{\Lambda}\min_{x,B}\
\mathcal{L}(x,B,\Lambda),\qquad
\mathcal{L}(x,B,\Lambda)
:=F(x,B)+\langle B,\Lambda\rangle-h^*(\Lambda).
\tag{D}
\]
We also define the critical set.
\(
\operatorname{crit}\mathcal{L}
:=\bigl\{(x,B,\Lambda):\
0\in\partial\mathcal{L}(x,B,\Lambda)\bigr\}.
% ,\quad\text{i.e.,}\quad
% \nabla_xF=0,\quad
% \nabla_BF+\Lambda=0,\quad
% B\in\partial h^*(\Lambda).
\)

Lemma 5.2 motivates our complexity analysis for finding a point
$(\hat x,\widehat B,\widehat\Lambda)$ satisfying 
\[
\mathbb{E}[\operatorname{dist}(0,\,\partial\mathcal{L}(\hat x,\widehat B,\widehat\Lambda))^2]\le\varepsilon^2.
\]
This constitutes the main body of our theoretical results in the submitted manuscript. 
Indeed, we agree that discussing whether a zero-gap primal--dual optimal pair
exists is important, as strong duality may fail due to the nonconvex nature of
the W-regularizer $h$ in Problem (P). We then illustrate this point with a
concrete instance.

\begin{example}\label{ex:gap}
Consider the one-dimensional instance
$p^* = \min_{b\in\mathbb{R}}\bigl\{\tfrac12 b^2+\bigl||b|-1\bigr|\bigr\}$,
i.e., $F(b)=\tfrac12 b^2$ and $\lambda=1$. The primal optimum is attained at
$b^*=\pm1$ with $p^*=\tfrac12$. Since
$\min_b\{\tfrac12 b^2+\Lambda b\}=-\tfrac12\Lambda^2$ and
$h^*(\Lambda)=|\Lambda|$ on $[-1,1]$,
\[
d^*=\sup_{|\Lambda|\le1}\Bigl\{-\tfrac12\Lambda^2-|\Lambda|\Bigr\}=0,
\]
attained at $\Lambda^*=0$. Hence $p^*-d^*=\tfrac12>0$: strong duality
between Problem~(P) and Problem~(D) fails, and no zero-gap primal--dual pair exists.
\end{example}

Nevertheless, the reviewer's question highlights an important point: what does $\operatorname{crit}\mathcal{L}$ imply for the 
original problem~\textup{(P)}? We appreciate this question, which prompted
us to establish the following explicit characterization. The proof is straightforward and closely follows the argument in Lemma~5.2 of the submitted manuscript.

\begin{proposition}\label{prop:transfer}
Let $(\tilde x,\widetilde B,\widetilde\Lambda)\in\operatorname{crit}\mathcal{L}$,
i.e.,
\begin{equation}\label{eq:crit}
\nabla_xF(\tilde x,\widetilde B)=0,\qquad
\nabla_BF(\tilde x,\widetilde B)+\widetilde\Lambda=0,\qquad
\widetilde B\in\partial h^*(\widetilde\Lambda).
\end{equation}
Then the following statements are equivalent:
\begin{enumerate}[label=(\alph*)]
\item $(\tilde x,\widetilde B)$ is a stationary point of Problem~(P), i.e.,
\begin{equation}\label{eq:pstat}
\nabla_xF(\tilde x,\widetilde B)=0,\qquad
0\in\nabla_BF(\tilde x,\widetilde B)+\partial h(\widetilde B);
\end{equation}
\item the Fenchel--Young equality in Lemma~5.2 holds at
$(\widetilde B,\widetilde\Lambda)$:
$h(\widetilde B)+h^*(\widetilde\Lambda)=\langle\widetilde\Lambda,\widetilde B\rangle$;
\item $|\widetilde B_{ij}|\ge1$ for all $i,j$.
\end{enumerate}
\end{proposition}
\begin{proof} 
Since $h^{**}$ is proper, closed and convex, and $(h^{**})^*=h^*$, the
subgradient inversion rule~\cite[Theorem~23.5]{rockafellar1970convex} gives
\[
\widetilde{B}\in\partial h^*(\widetilde{\Lambda})
\iff
\widetilde{\Lambda}\in\partial h^{**}(\widetilde{B})
\iff
h^{**}(\widetilde{B})+h^*(\widetilde{\Lambda})
=\langle\widetilde{\Lambda},\widetilde{B}\rangle.
\tag{3}
\]
In addition, \eqref{eq:crit} gives
$\widetilde{\Lambda}=-\nabla_B F(\tilde{x},\widetilde{B})$, so
\eqref{eq:pstat} holds if and only if
$\widetilde{\Lambda}\in\partial h(\widetilde{B})$.

\emph{(b)$\iff$(c).} By (3), the Fenchel--Young equality for $h$ holds at
$(\widetilde{B},\widetilde{\Lambda})$ if and only if
$h(\widetilde{B})=h^{**}(\widetilde{B})$. Since the functions are separable,
comparing $h$ with $h^{**}(B)=\lambda\sum_{i,j}(|B_{ij}|-1)_+$ componentwise
shows that this is equivalent to $|\widetilde{B}_{ij}|\ge1$ for all $i,j$.

\emph{(b)$\Rightarrow$(a).} The equality
$h(\widetilde{B})+h^*(\widetilde{\Lambda})
=\langle\widetilde{\Lambda},\widetilde{B}\rangle$ together with the
definition of $h^*$ means that $\widetilde{B}$ globally minimizes
$h(\cdot)-\langle\widetilde{\Lambda},\cdot\rangle$. By Fermat's
rule~\cite[Theorem~10.1]{rockafellar1998variational} and Proposition~A.2 of
the manuscript, $\widetilde{\Lambda}\in\widehat{\partial}h(\widetilde{B})
\subseteq\partial h(\widetilde{B})$, which is \eqref{eq:pstat}.

\emph{(a)$\Rightarrow$(c).}  (a) means
$\widetilde{\Lambda}\in\partial h(\widetilde{B})
\cap\partial h^{**}(\widetilde{B})$. Both functions are separable, so it
suffices to examine a scalar component $b=\widetilde{B}_{ij}$. If $|b|<1$,
then $\partial h^{**}(b)=\{0\}$, whereas
$\partial h(b)=\{-\lambda\,\mathrm{sgn}(b)\}$ for $b\neq0$ and
$\partial h(0)=\{-\lambda,\lambda\}$; in all cases $0\notin\partial h(b)$,
so the intersection is empty. Hence $|\widetilde{B}_{ij}|\ge1$ for all $i,j$.
\end{proof}

We highlight that Proposition~\ref{prop:transfer}, though quite simple, may answer the reviewer's question: it shows that a critical point $(\tilde x,\widetilde B,\widetilde\Lambda)$ of $\mathcal{L}$ yields a first-order stationary point $(\tilde x,\widetilde B)$ of Problem~(P) if and
only if $|\widetilde B_{ij}|\ge1$ for all $i,j$. Moreover, our numerical experiments indicate that the primal--dual formulation is practically
advantageous, yielding lower quantization errors and higher retrieval
precision, which suggests that optimizing in the dual space is more effective than directly optimizing the nonconvex W-regularizer in Problem (P).

In summary, we appreciate the reviewer’s comment. However, we believe the submitted manuscript has already partially addressed this question. We are also prepared to add Example 1 and Proposition 1 to the revised manuscript. Nevertheless, our main theoretical contributions regarding computational complexity and numerical performance of DualHash remain unaffected. 


\section{Reviewer 2: Comments on Assumptions and Corrections to the Theoretical Analysis}

\vspace{0.5em}
\noindent\textbf{Main comment.} \emph{%
The convergence and complexity analysis has several issues, including
insufficient assumptions, incomplete boundedness and
descent arguments, and some errors in conditional expectations.
Some of these affect the main complexity results, so the
theoretical guarantees are not fully established as written.
I return the article to the authors for \textbf{major revisions}.}

\vspace{0.5em}
\noindent\textbullet\quad\textbf{Comment 1.}

\noindent\emph{\textbf{Coercivity and lower-boundedness.}
(Lemma~5.1 and Lemmas~5.7, 5.10, D.2)
The boundedness argument is incomplete: the coercivity and
lower-boundedness of $J$ alone do not imply the boundedness of
$\{(x^k,B^k)\}$ without controlling
$\{J(x^k,B^k)\}$, and coercivity is not included in
Assumption~5.1. In addition, the lower bounds required to define
$\Psi_m^*$ and $\Psi_r^*$, as well as the non-increasing estimate for
$\{\mathbb{E}[\Psi_r(\mathcal{Q}^k)]\}$ in Lemma~D.2, require
justification.}

% \begin{enumerate}[label=(\arabic*),leftmargin=*,labelsep=0.5em,topsep=0.25em,itemsep=0.25em]
% \item \textbf{Boundedness.} \emph{(Lemma 5.1/proof) Having $J$ coercive and lower-bounded is generally not
% enough to directly obtain the boundedness of $(x^k, B^k)$. You classically also need
% to prove that the sequence $\{J(x^k, B^k)\}$ converges (or at least its boundedness). Note
% also that having boundedness of all level-sets is equivalent to coercivity. Finally,
% the coercivity of $J$ is not assumed in your Assumption~5.1.}

% \item \textbf{Lower-boundedness.} \emph{(Lemmas 5.7, 5.10, D.2) The lower-boundedness of
% $\{\Psi_m(\mathcal{Q}^k)\}$ and $\{\Psi_r(\mathcal{Q}^k)\}$ needs to be
% justified to invoke $\Psi_m^*$ and $\Psi_r^*$, which doesn't seem
% straightforward. Similarly, the non-increasing property of
% $\{\mathbb{E}[\Psi_r(\mathcal{Q}^k)]\}$ in Lemma~D.2 requires
% justification.}
% \end{enumerate}

\vspace{0.5em}
\noindent\textbullet\quad\textbf{Our response.}
We thank the reviewer for the important comment. While we agree that deducing boundedness of the iterates in Lemma~5.1 from coercivity of the objective function $J(x, B)$ requires a uniform upper bound on the objective values along the iterates, our complexity analysis actually does not require the iterates themselves to be bounded. Rather, it just requires lower bounds on the potentials used in the analysis. This was the original motivation for introducing Lemma~5.1.

Regarding coercivity, we already stated in the submitted manuscript
(p.~7) that the objective $J$ in~(3.1),
\[
J(x,B)=
\underbrace{\frac1n\sum_{i=1}^n\ell_{s,i}(x)}_{=:\,\bar{\ell}_s(x)}
+\frac{\mu}{2}\|x\|^2
+\frac{\gamma}{2n}\sum_{i=1}^n\|\mathcal D_i(x)-b_i\|^2
+h(B),
\]
is coercive. However, to ensure this property through the model
structure, we will impose the lower-boundedness assumption on
the average loss rather than directly on the objective:
\[
\underline{\ell}:=\inf_x\bar{\ell}_s(x)>-\infty.
\]
Together with the model structure, this condition ensures both
coercivity and lower-boundedness of $J$. We will therefore refine
Assumption~5.1(i) accordingly. This refinement is intended to
rule out cancellation of the growth of the regularization terms
by the loss and to justify the coercivity already stated in
the manuscript.

Under the refined Assumption~5.1(i) and the
appropriate parameter conditions, the estimates in our analysis
yield the required lower bounds on the potentials. We will make
these estimates explicit in the revision. Nevertheless, these
adjustments preserve the stated complexity orders and the main
theoretical contributions of the submitted manuscript.
% \noindent\textbf{(i) Coercivity of $J$.}

% As stated in the submitted manuscript (p.~7), the model
% structure~(3.1),
% \[
% \min_{x,B}\;
% J(x,B)=F(x,B)+h(B),\quad
% F(x,B)=\frac{1}{n}{\textstyle\sum_{i=1}^{n}}f_i(x)
% +\frac{\gamma}{2n}{\textstyle\sum_{i=1}^{n}}\|\mathcal{D}_i(x)-b_i\|^2,
% \]
% where $f_i(x)=\ell_{s,i}(x)+\frac{\mu}{2}\|x\|^2$ with $\mu>0$ and
% $\bar\ell_s(x):=\frac{1}{n}\sum_{i=1}^{n}\ell_{s,i}(x)$, already entails the coercivity of~$J$ under the following
% refined assumption, which replaces the lower-boundedness of~$J$ in
% Assumption~5.1(i) (p.~10):
% \[
% \underline{\ell}\;:=\;\inf_{x}\,\bar\ell_s(x)\;>\;-\infty.
% \]
% Indeed, since the W-type regularizer satisfies
% $h(B)\ge\lambda\|B\|_1-\lambda nd$, we have
% \[
% J(x,B)\;\ge\;\underline{\ell}
% +\frac{\mu}{2}\|x\|^2+\lambda\|B\|_1-\lambda nd
% \;\longrightarrow\;+\infty
% \quad\text{as }\|(x,B)\|\to\infty,
% \]
% so the coercivity and lower-boundedness of~$J$ need not be
% assumed separately.

% \smallskip
% \noindent\textbf{(ii) Lower-boundedness of the Lyapunov functions.}
% We now provide a complete justification of the lower-boundedness of the Lyapunov functions $\Psi_m(\mathcal{Q}^k)$ in~(5.12) and $\Psi_r(\mathcal{Q}^k)$ in~(5.21) under the revised assumptions, as follows.

% \begin{proposition}
% \label{prop:lyapunov-lower-bound}
% Under the revised Assumption~5.1 (p.~10) and
% step-size conditions compatible with those in Theorem~5.2 (p.~14) and
% Theorem~5.3 (p.~16),
% \begin{enumerate}[label=\textup{(\roman*)},nosep]
% \item for DualHash-StoM, the following lower bound holds for every $k\ge 1$:
% \[
% \mathbb{E}[\Psi_m(\mathcal{Q}^k)]
% \;\ge\;
% \underline{\ell}
% +\frac{\mu}{4}\,\mathbb{E}\|x^k\|^2
% +\frac{3\gamma}{16n}\,\mathbb{E}\|B^k
%   -\mathcal{D}(x^k)\|_F^2
% -c_1
% \;\ge\;
% \Psi_m^*>-\infty.
% \]
% \item for DualHash-StoRM, let $e^k:=D^k-\nabla_xF(x^k,B^k)$ denote the
% estimation error of the StoRM estimator, and let $\eta_T$ and $\rho_T$ be the
% step size and the momentum parameter of the schedule in Theorem~5.3 (p.~16).
% Then the
% combined potential
% $W_k:=\mathbb{E}[\Psi_r(\mathcal{Q}^k)]
% +\frac{\eta_T}{\rho_T}\,\mathbb{E}\|e^k\|^2$
% satisfies
% \[
% W_k\;\ge\;\Psi_r^*>-\infty,
% \qquad k\ge 1.
% \]
% \end{enumerate}
% The constants $c_1$, $\Psi_m^*$, and $\Psi_r^*$ are independent
% of~$T$. 
% \end{proposition}

\vspace{0.5em}
\noindent\textbullet\quad\textbf{Comments 2 and 3
(Descent argument; conditional expectations and estimates).}

\noindent\emph{%
The reviewer pointed out several issues with the proofs, including
an incorrect summation formula in Lemma~5.7 that
propagates into Theorem~5.2; insufficient justification
that $C_B>0$ ensures $C_2>0$ (Lemma~C.2);
misuse of conditional vs.\ full expectations in
Appendices~C and~D; and the need to rework the proofs
of Lemma~5.10 and Theorem~5.3 based on the corrected
inequality~(D.4).}

\vspace{0.5em}
\noindent\textbullet\quad\textbf{Our response.}
We thank the reviewer for the careful review. We will correct
the estimates for the expectations and variances and refine the
initialization and parameter conditions accordingly.
These corrections include replacing the conditional expectations
in~(C.7) with full expectations, using the same batch for both
sampled gradients in~(4.8), and restoring the omitted term in
the variance-reduction recursion.
Nevertheless, these adjustments preserve the
$\mathcal{O}(\varepsilon^{-3})$ SFO complexity stated in
Theorem~5.3. We will incorporate these corrections into the revised manuscript.

% We thank the reviewer for identifying these errors; the corrections are
% confined to constants and conditional-expectation formalism and do not affect the established complexity orders, as detailed below.

% \smallskip
% \noindent\textbf{(i) Coefficient and step-size condition (Lemma~C.2, p.~32).}
% The corrected coefficient and strict step-size condition read
% \[
% \frac{2C_B}{L_F}=\blue{\frac1t}-\frac{13}{2}t-2>0
% \quad\text{when}\quad
% 0<\tau\blue{<}\frac{\sqrt{30}-2}{13L_F},
% \qquad t:=\tau L_F.
% \]
% The corrected averaged estimate retains the same form as the estimate stated in Theorem~5.2 (p.~14),
% with
% updated constants $L_\Delta$ and $C_\sigma$; for $b_k=c_bT$, the
% $\mathcal{O}(\varepsilon^{-4})$ SFO complexity order is still preserved.

% \smallskip
% \noindent\textbf{(ii) Conditional expectations and batch index.}
% The conditional expectations in~(C.7) are replaced by full expectations~$\mathbb{E}[\,\cdot\,]$. In addition, the corrected StoRM estimator in~(4.8) reads
% \[
% D^k=\begin{cases}
% \nabla_x F(x^1,B^1;\mathcal{J}_1), & k=1,\\[4pt]
% (1-\rho_{k-1})\bigl(D^{k-1}
%   -\nabla_x F(x^{k-1},B^{k-1};\blue{\mathcal{J}_k})\bigr)
%   +\nabla_x F(x^k,B^k;\mathcal{J}_k), & k>1.
% \end{cases}
% \]

% \smallskip
% \noindent\textbf{(iii) Variance-reduction recursion.} The
% corrected Lemma~D.1 (p.~36) restores the $\blue{\Delta_B^{k+1}}$ term omitted
% in the variance-reduction recursion. With $m:=\gamma/n$, the explicit
% initialization
% \[
% \blue{B^1=\operatorname{prox}_{m^{-1}h^{**}}(\mathcal{D}(x^1))},
% \qquad
% \blue{\Lambda^1=m(\mathcal{D}(x^1)-B^1)},
% \]
% and carefully refined parameter conditions, we still obtain
% \[
% \frac1T\sum_{k=1}^{T}\mathbb{E}\operatorname{dist}^2
% \bigl(0,\partial\mathcal{L}(x^k,B^k,\Lambda^k)\bigr)
% \le\frac{C}{T^{2/3}},
% \]
% where $C$ is independent of~$T$. The $\mathcal{O}(\varepsilon^{-3})$ SFO
% complexity in Theorem~5.3 (p.~16) is therefore preserved.

\bigskip

%\section*{}
In closing, deep hashing maps images to short binary codes (e.g., 64 bits) for large-scale retrieval, reducing storage by several orders of magnitude and enabling
sub-millisecond search via Hamming distance. The W-type regularizer
$\|\,|z|-\mathbf{1}\|_1$ is the principal regularizer used in practice to enforce
binarization, yet its nonconvexity and nonsmoothness have left convergence and
complexity analysis largely unexplored for W-regularized hashing methods \cite{liu2016dsh,zhu2016dhn,wang2023mds,chen2025shc}. Our manuscript addresses
this gap: we reformulate the problem via Fenchel duality as a convex-dual formulation with
closed-form proximal updates, and establish the first complexity guarantees for
this setting, complemented by extensive experiments whose practical value both
reviewers acknowledged.

We would be grateful for the opportunity to submit a revised manuscript. 

\vspace{1.5em}

\noindent Sincerely,\\
Authors of \textbf{Submission \#M187198}.

\bibliographystyle{siamplain}
\bibliography{revision_refs}

% \section*{Appendix 1: Proof of Proposition~\ref{prop:transfer}}

% \begin{proof}
% \emph{Step 0: a consequence of \eqref{eq:crit}.}
% Since $h^{**}$ is proper, closed and convex, and $(h^{**})^*=h^*$, the
% subgradient inversion rule~\cite[Theorem~23.5]{rockafellar1970convex} gives
% \[
% \widetilde{B}\in\partial h^*(\widetilde{\Lambda})
% \iff
% \widetilde{\Lambda}\in\partial h^{**}(\widetilde{B})
% \iff
% h^{**}(\widetilde{B})+h^*(\widetilde{\Lambda})
% =\langle\widetilde{\Lambda},\widetilde{B}\rangle.
% \tag{A.1}
% \]
% In addition, \eqref{eq:crit} gives
% $\widetilde{\Lambda}=-\nabla_B F(\tilde{x},\widetilde{B})$, so
% \eqref{eq:pstat} holds if and only if
% $\widetilde{\Lambda}\in\partial h(\widetilde{B})$.

% \emph{(b)$\iff$(c).} By (A.1), the Fenchel--Young equality for $h$ holds at
% $(\widetilde{B},\widetilde{\Lambda})$ if and only if
% $h(\widetilde{B})=h^{**}(\widetilde{B})$. Comparing $h$ with
% $h^{**}(B)=\lambda\sum_{i,j}(|B_{ij}|-1)_+$ componentwise, this is equivalent
% to $|\widetilde{B}_{ij}|\ge1$ for all $i,j$.

% \emph{(b)$\Rightarrow$(a).} The equality
% $h(\widetilde{B})+h^*(\widetilde{\Lambda})
% =\langle\widetilde{\Lambda},\widetilde{B}\rangle$ together with the
% definition of $h^*$ means that $\widetilde{B}$ globally minimizes
% $h(\cdot)-\langle\widetilde{\Lambda},\cdot\rangle$. By Fermat's
% rule~\cite[Theorem~10.1]{rockafellar1998variational} and Proposition~A.2
% (p.~29) of the manuscript, $\widetilde{\Lambda}\in\widehat{\partial}h(\widetilde{B})
% \subseteq\partial h(\widetilde{B})$, which is \eqref{eq:pstat}.

% \emph{(a)$\Rightarrow$(c).} By Step~0, (a) means
% $\widetilde{\Lambda}\in\partial h(\widetilde{B})
% \cap\partial h^{**}(\widetilde{B})$. Both functions are separable, so it
% suffices to examine a scalar component $b=\widetilde{B}_{ij}$. If $|b|<1$,
% then $\partial h^{**}(b)=\{0\}$, whereas
% $\partial h(b)=\{-\lambda\,\mathrm{sgn}(b)\}$ for $b\neq0$ and
% $\partial h(0)=\{-\lambda,\lambda\}$; in all cases $0\notin\partial h(b)$,
% so the intersection is empty. Hence $|\widetilde{B}_{ij}|\ge1$ for all $i,j$.
% \end{proof}

% \clearpage
% \section*{Appendix 2: Proof of Proposition~\ref{prop:lyapunov-lower-bound}}

% Let $m:=\gamma/n$, $M:=\lambda\sqrt{nd}$, and
% $C_D:=\sup_x\|\mathcal{D}(x)\|_F$. The uniform output condition in the
% revised Assumption~B.1 (p.~30) gives $C_D<\infty$; for the hashing architecture
% with a final $\tanh$ activation, one may take $C_D\le\sqrt{nd}$. On
% $\operatorname{dom}h^*=[-\lambda,\lambda]^{n\times d}$, we have
% \[
% \|\Lambda\|_F\le M,
% \qquad
% 0\le h^*(\Lambda)\le\lambda nd.
% \]

% \subsection*{2.1 DualHash-StoM}

% \begin{proof}[Proof of Proposition~\ref{prop:lyapunov-lower-bound}(i)]
% Write
% \[
% r_k:=B^k-\mathcal{D}(x^k),
% \qquad
% p_k:=x^k-x^{k-1},
% \qquad
% u_k:=x^{k+1}-x^k.
% \]
% Recall the Lyapunov function~(5.12) on page~13 of the submitted manuscript,
% \[
% \Psi_m(\mathcal{Q}^k)=\mathcal{L}(x^k,B^k,\Lambda^k)
% -C_1\Delta_B^{k+1}+C_2\Delta_B^k
% -C_3\Delta_x^{k+1}+C_4\Delta_x^k.
% \]

% \emph{Step 1: a lower bound for the Lagrangian.}
% Since $\|B^k\|_F\le\|r_k\|_F+C_D$, Young's inequality gives
% \[
% \langle B^k,\Lambda^k\rangle-h^*(\Lambda^k)
% \ge-\frac m8\|r_k\|_F^2-c_0,
% \qquad
% c_0:=\frac{2M^2}{m}+MC_D+\lambda nd.
% \]
% Moreover,
% \[
% F(x^k,B^k)\ge\underline\ell+\frac\mu2\|x^k\|^2
% +\frac m2\|r_k\|_F^2.
% \]
% Hence
% \begin{equation}\label{eq:prop2-L-lower}
% \mathcal{L}(x^k,B^k,\Lambda^k)
% \ge\underline\ell+\frac\mu2\|x^k\|^2
% +\frac{3m}{8}\|r_k\|_F^2-c_0.
% \end{equation}

% \emph{Step 2: the forward $B$-increment.}
% The $B$-update and the Lipschitz continuity of $\mathcal{D}$ give
% \[
% \begin{aligned}
% \Delta_B^{k+1}
% &=\frac{\tau^2}{2}
% \|m(B^k-\mathcal{D}(x^{k+1}))+\Lambda^k\|_F^2\\
% &\le\frac{3\tau^2}{2}
% \left(m^2\|r_k\|_F^2+m^2L_D^2\|u_k\|^2+\lambda^2nd\right).
% \end{aligned}
% \]
% Define
% \[
% A_B:=\frac{3C_1\tau^2m^2}{2},
% \qquad
% H_m:=\frac{3C_1\tau^2m^2L_D^2}{2}+\frac{C_3}{2}.
% \]
% Then
% \begin{equation}\label{eq:prop2-B-increment}
% -C_1\Delta_B^{k+1}-C_3\Delta_x^{k+1}
% \ge-A_B\|r_k\|_F^2-H_m\|u_k\|^2
% -\frac{3C_1\tau^2\lambda^2nd}{2}.
% \end{equation}

% \emph{Step 3: the stochastic momentum increment.}
% Let
% $z^k=x^k+\beta p_k$ and
% $e^k=G^k-\nabla_xF(z^k,B^k)$. The $x$-update gives
% \[
% u_k=\alpha p_k-\eta_k\bigl(\nabla_xF(z^k,B^k)+e^k\bigr).
% \]
% By global $L_F$-smoothness and $F\ge\underline\ell$,
% \[
% \|\nabla F(x^k,B^k)\|^2
% \le2L_F\bigl(F(x^k,B^k)-\underline\ell\bigr),
% \]
% and therefore
% \[
% \begin{aligned}
% \mathbb{E}\|u_k\|^2
% \le{}&12\eta_k^2L_F\,
% \mathbb{E}\bigl[F(x^k,B^k)-\underline\ell\bigr]\\
% &+\bigl(3\alpha^2+6\eta_k^2L_F^2\beta^2\bigr)
% \mathbb{E}\|p_k\|^2
% +3\eta_k^2\frac{\sigma^2}{b_k}.
% \end{aligned}
% \]
% Put
% \[
% \Theta_k:=12H_m\eta_k^2L_F,
% \qquad
% Q_k:=3\alpha^2+6\eta_k^2L_F^2\beta^2.
% \]
% We impose the following explicit absorption conditions:
% \begin{align}
% \Theta_k&\le\frac12, \tag{P1}\label{eq:prop2-P1}\\
% \frac m2(1-\Theta_k)-\frac m8-A_B&\ge\frac{3m}{16},
% \tag{P2}\label{eq:prop2-P2}\\
% C_4&\ge2H_mQ_k. \tag{P3}\label{eq:prop2-P3}
% \end{align}
% They are compatible with the stated parameter restrictions after the fixed
% constants are chosen sufficiently small.

% Combining \eqref{eq:prop2-L-lower}, \eqref{eq:prop2-B-increment}, and
% the preceding estimate yields
% \[
% \begin{aligned}
% \mathbb{E}\Psi_m(\mathcal{Q}^k)
% \ge{}&\underline\ell+\frac\mu4\mathbb{E}\|x^k\|^2
% +\frac{3m}{16}\mathbb{E}\|r_k\|_F^2
% +C_2\mathbb{E}\Delta_B^k+\frac{C_4}{2}\mathbb{E}\Delta_x^k\\
% &-c_0-\frac{3C_1\tau^2\lambda^2nd}{2}
% -3H_m\eta_k^2\frac{\sigma^2}{b_k}.
% \end{aligned}
% \]
% Since $b_k\ge1$, the final term is bounded below by a constant independent
% of $T$. Dropping the nonnegative terms proves
% \(\mathbb{E}\Psi_m(\mathcal{Q}^k)\ge\Psi_m^*>-\infty\), where
% \(\Psi_m^*\) is independent of $T$.
% \end{proof}

% \subsection*{2.2 DualHash-StoRM}

% \begin{proof}[Proof of Proposition~\ref{prop:lyapunov-lower-bound}(ii)]
% Recall the Lyapunov function~(5.21) on page~15 of the submitted manuscript,
% \[
% \Psi_r(\mathcal{Q}^k)=\mathcal{L}(x^k,B^k,\Lambda^k)
% -C_1\Delta_B^{k+1}+C_2\Delta_B^k-C_3\Delta_x^{k+1}
% \]
% and define
% \[
% W_k:=\mathbb{E}[\Psi_r(\mathcal{Q}^k)]
% +\frac{\eta_T}{\rho_T}\mathbb{E}\|e^k\|^2,
% \qquad
% e^k:=D^k-\nabla_xF(x^k,B^k).
% \]
% The bounds \eqref{eq:prop2-L-lower} and the forward $B$-increment estimate
% remain valid. Since $u_k=-\eta_kD^k$,
% \[
% \mathbb{E}\|u_k\|^2
% \le2\eta_k^2\left(
% \mathbb{E}\|\nabla_xF(x^k,B^k)\|^2+\mathbb{E}\|e^k\|^2\right).
% \]
% Set
% \[
% \Gamma_k:=\eta_k^2\bigl(3C_1\tau^2m^2L_D^2+C_3\bigr).
% \]
% Using the same notation $A_B=3C_1\tau^2m^2/2$ as above, we obtain
% \[
% \begin{aligned}
% W_k\ge{}&\mathbb{E}\bigl[F(x^k,B^k)\bigr]
% -\left(\frac m8+A_B\right)\mathbb{E}\|r_k\|_F^2-c_0
% -\frac{3C_1\tau^2\lambda^2nd}{2}\\
% &-\Gamma_k\mathbb{E}\|\nabla_xF(x^k,B^k)\|^2
% +\left(\frac{\eta_T}{\rho_T}-\Gamma_k\right)\mathbb{E}\|e^k\|^2
% +C_2\mathbb{E}\Delta_B^k.
% \end{aligned}
% \]
% The following conditions are sufficient:
% \begin{align}
% \Gamma_k&\le\frac{\eta_T}{\rho_T},
% \tag{P1$'$}\label{eq:prop2-P1storm}\\
% 2L_F\Gamma_k&\le\frac12,
% \tag{P2$'$}\label{eq:prop2-P2storm}\\
% \frac m2(1-2L_F\Gamma_k)-\frac m8-A_B&\ge\frac{3m}{16}.
% \tag{P3$'$}\label{eq:prop2-P3storm}
% \end{align}
% Indeed, the global smoothness bound
% \[
% \|\nabla_xF(x^k,B^k)\|^2
% \le2L_F\bigl(F(x^k,B^k)-\underline\ell\bigr)
% \]
% allows the gradient-square term to be absorbed by the positive terms in
% \eqref{eq:prop2-L-lower}, while \emph{(P1$'$)} absorbs the error-square
% term into the tracking term in $W_k$. Under the parameter schedule of
% Theorem~5.3, $\eta_k=O(T^{-1/3})$ and
% $\eta_T/\rho_T=O(T^{1/3})$; hence these conditions are compatible with
% the stated schedule after the fixed constants are chosen sufficiently small.

% It follows that
% \[
% W_k\ge\underline\ell+\frac\mu4\mathbb{E}\|x^k\|^2
% +\frac{3m}{16}\mathbb{E}\|B^k-\mathcal{D}(x^k)\|_F^2
% +C_2\mathbb{E}\Delta_B^k-c_1',
% \]
% where
% \[
% c_1':=c_0+\frac{3C_1\tau^2\lambda^2nd}{2}
% \]
% is independent of $T$. Dropping the nonnegative terms proves
% \(W_k\ge\Psi_r^*:=\underline\ell-c_1'>-\infty\).
% \end{proof}

% % \clearpage
% \section*{Appendix B: Detailed Descent Argument and Corrections}
% \noindent\emph{Internal working notes for the manuscript revision; not part
% of the response to the reviewer.}

% \smallskip
% \noindent\textbf{(i) Intermediate steps leading to~(5.9).}
% Recall that
% \[
% \mathcal{L}(x,B,\Lambda)
% =F(x,B)+\langle B,\Lambda\rangle-h^*(\Lambda).
% \]
% The optimality condition for the preceding dual update gives
% \[
% s^k:=2B^k-B^{k-1}+\tau(\Lambda^{k-1}-\Lambda^k)
% \in\partial h^*(\Lambda^k),
% \]
% Since $h^*$ is convex, its subgradient inequality yields
% \[
% -h^*(\Lambda^{k+1})+h^*(\Lambda^k)
% \le -\langle s^k,\Lambda^{k+1}-\Lambda^k\rangle.
% \]
% Applying the descent lemma to $F(x^{k+1},\cdot)$, and then using
% \[
% B^{k+1}-B^k
% =-\tau\bigl(\nabla_BF(x^{k+1},B^k)+\Lambda^k\bigr),
% \]
% we obtain
% \[
% \begin{aligned}
% &F(x^{k+1},B^{k+1})-F(x^{k+1},B^k)
% +\langle\Lambda^k,B^{k+1}-B^k\rangle\\
% &\qquad\le
% -\left(\tau^{-1}-\frac{L_F}{2}\right)
% \|B^{k+1}-B^k\|^2.
% \end{aligned}
% \]
% Combining these two inequalities gives
% \[
% \begin{aligned}
% &\mathcal{L}(x^{k+1},B^{k+1},\Lambda^{k+1})
% -\mathcal{L}(x^{k+1},B^k,\Lambda^k)\\
% &\quad\leq
% -\left(\tau^{-1}-\frac{L_F}{2}\right)
% \|B^{k+1}-B^k\|^2
% +\langle\Lambda^{k+1}-\Lambda^k,B^{k+1}-s^k\rangle.
% \end{aligned}
% \]
% The remaining term follows from two consecutive $B$-updates:
% \[
% \begin{aligned}
% B^{k+1}-s^k
% &=(B^{k+1}-B^k)-(B^k-B^{k-1})
% +\tau(\Lambda^k-\Lambda^{k-1})\\
% &=\tau\bigl(
% \nabla_BF(x^k,B^{k-1})
% -\nabla_BF(x^{k+1},B^k)
% \bigr).
% \end{aligned}
% \]
% Substitution gives exactly~(5.9). We will include these steps in the revised
% proof and make the starting index consistent with the initialization, since
% this argument uses the preceding primal and dual updates.

% \smallskip
% \noindent\textbf{(ii) Positivity argument in Lemma~C.2.}
% With $t=\tau L_F$, the definitions of $K_1,K_2,K_3$ give
% \[
% C_B
% =\frac{K_2-K_3-K_1}{2}
% =\frac{L_F}{2}
% \left[
% \frac{2-6\delta}{t}
% -\left(3\delta+\frac1\delta\right)t
% -(6\delta+1)
% \right].
% \]
% For $\delta=1/6$, the correct expression is
% \[
% \frac{2C_B}{L_F}
% =\frac1t-\frac{13}{2}t-2.
% \]
% Consequently, $C_B>0$ holds whenever
% \[
% 0<\tau<\frac{\sqrt{30}-2}{13L_F}.
% \]
% The upper endpoint must be excluded to ensure strict positivity. Furthermore,
% \[
% C_2=C_B+K_3,
% \qquad
% K_3=\frac{\tau L_F^2}{\delta}>0,
% \]
% which immediately establishes $C_2>0$. We will correct the displayed
% coefficient, remove the premature assertions of nonnegativity, and state an
% explicit admissible parameter range in the revised manuscript.

% \smallskip
% \noindent\textbf{(iii) Noise terms and their accumulation in Lemma~5.7.}
% We agree that, for varying batch sizes, the averaged estimate must retain
% \[
% \frac1T\sum_{k=1}^{T}\frac1{b_k}.
% \]
% Tracing this term back to Lemma~C.3 also reveals an inconsistency between
% its statement and proof. The proof contains $3\|e^k\|^2$, whose expectation
% is bounded by $3\sigma^2/b_k$. Thus, the factor $L_F^{-1}$ in the
% corresponding noise term of the lemma statement should be removed.
% Combining the corrected stationarity estimate with the Lyapunov descent
% inequality gives
% \[
% \begin{aligned}
% &\frac1T\sum_{k=1}^{T}
% \mathbb{E}\!\left[
% \operatorname{dist}^2
% \bigl(0,\partial\mathcal{L}(x^k,B^k,\Lambda^k)\bigr)
% \right]\\
% &\qquad\leq
% \frac{L_\Delta\Delta_1}{T}
% +\left(3+\frac{L_\Delta}{L_F}\right)
% \frac{\sigma^2}{T}\sum_{k=1}^{T}\frac1{b_k},
% \end{aligned}
% \]
% provided that the required lower bound is established in the separate boundedness argument.
% We will also align the averaging indices and the random output selection
% throughout the statement and proof.

% For the constant batch size $b_k=b=c_bT$ used in Theorem~5.2, this becomes
% \[
% \frac{L_\Delta\Delta_1}{T}
% +\left(3+\frac{L_\Delta}{L_F}\right)
% \frac{\sigma^2}{c_bT}.
% \]
% Therefore, the summation correction preserves the claimed iteration and SFO
% complexity orders in Theorem~5.2. We will revise the affected statements and
% proofs consistently.

% \clearpage
% \section*{Appendix C: Stochastic Estimates and a Sufficient Rate-Preserving Repair}
% \noindent\emph{Internal working notes for the manuscript revision. The batch
% index in (4.8) has been confirmed by the authors: both gradient evaluations
% use the current batch. The rate below concerns an explicitly revised parameter
% schedule; it does not establish the submitted fixed-$\tau$ theorem as written.}

% \smallskip
% \noindent\textbf{C.1. Information and expectation conventions.}
% Let $\mathcal{F}_k$ contain the initialization and all batches through
% iteration $k$. The iterate $(x^k,B^k,\Lambda^k)$ is
% $\mathcal{F}_{k-1}$-measurable, and the next batch is sampled independently
% of this history. Within each batch, use independent uniform samples with
% replacement. Assume all expectations below are finite. For StoM,
% \[
% \mathbb{E}[e^k\mid\mathcal{F}_{k-1}]=0,
% \qquad
% \mathbb{E}[\|e^k\|^2\mid\mathcal{F}_{k-1}]\leq\sigma^2/b_k.
% \]
% The Lyapunov variable $\Psi_m(\mathcal{Q}^k)$ contains next-iterate terms,
% so it must not be treated as $\mathcal{F}_{k-1}$-measurable. Take full
% expectations of its pathwise descent estimate to obtain the corrected (C.7):
% \[
% \begin{aligned}
% \mathbb{E}\Psi_m(\mathcal{Q}^{k+1})-\mathbb{E}\Psi_m(\mathcal{Q}^k)
% \leq{}&-C_B\mathbb{E}[\Delta_B^{k+1}+\Delta_B^k]\\
% &-C_x\mathbb{E}[\Delta_x^{k+1}+\Delta_x^k]
% +\frac{\sigma^2}{b_kL_F}.
% \end{aligned}
% \]
% No independence between the error and the next iterate is used. This
% restores the estimate needed for Theorem~5.2, subject to the separate
% lower-bound and initialization requirements.

% \smallskip
% \noindent\textbf{C.2. Same-batch StoRM recursion and the missing block.}
% Write $z^k=(x^k,B^k)$, $g(z)=\nabla_xF(x,B)$ and
% $g(z;\mathcal{J})=|\mathcal{J}|^{-1}\sum_{i\in\mathcal{J}}
% \nabla_xF_i(x,b_i)$. The intended update is
% \[
% D^{k+1}=g(z^{k+1};\mathcal{J}_{k+1})
% +(1-\rho_k)\bigl[D^k-g(z^k;\mathcal{J}_{k+1})\bigr].
% \]
% Here $z^{k+1},z^k,D^k$ are $\mathcal{F}_k$-measurable. Define
% $\epsilon(z;\mathcal{J})=g(z;\mathcal{J})-g(z)$. Then
% \[
% e^{k+1}=(1-\rho_k)e^k+\rho_k\epsilon(z^{k+1};\mathcal{J}_{k+1})
% +(1-\rho_k)\bigl[\epsilon(z^{k+1};\mathcal{J}_{k+1})
% -\epsilon(z^k;\mathcal{J}_{k+1})\bigr].
% \]
% The last two terms have conditional mean zero. Independence within the
% batch, the variance bound and joint mean-square smoothness imply
% \[
% \begin{aligned}
% \mathbb{E}[\|e^{k+1}\|^2\mid\mathcal{F}_k]
% \leq{}&(1-\rho_k)^2\|e^k\|^2+\frac{2\rho_k^2\sigma^2}{b_{k+1}}\\
% &+\frac{4(1-\rho_k)^2L_F^2}{b_{k+1}}
% \bigl(\Delta_x^{k+1}+\Delta_B^{k+1}\bigr).
% \end{aligned}
% \tag{C.1}
% \]
% Indeed, the variance of the centered gradient difference is bounded by
% $L_F^2\|z^{k+1}-z^k\|^2/b_{k+1}$; applying
% $\|u+v\|^2\leq2\|u\|^2+2\|v\|^2$ gives the displayed constants.
% Taking full expectations gives the corrected Lemma~D.1. The
% $\Delta_B^{k+1}$ term cannot be omitted under the stated joint smoothness
% assumption. For example, with $n=2$, scalar $x$, and
% $D_1(x)=D_2(x)=x$, changing only $b_1$ by $d$ at a fixed $x$ changes the
% centered single-sample $x$-gradient by $\pm\gamma d/2$. Its variance is
% $\gamma^2d^2/4>0$ even though the $x$-increment is zero. This example tests
% the claimed smoothness-based omission; it is not a counterexample to the
% full algorithm's convergence.

% \smallskip
% \noindent\textbf{C.3. The algebraic repair of (D.7).}
% If an $x$-only error recursion had been justified, the disputed step could
% be avoided by first weakening its nonnegative increment coefficient:
% \[
% q_{k+1}\leq(1-\rho_k)q_k
% +\frac{2\rho_k^2\sigma^2}{b_{k+1}}
% +4L_F^2\mathbb{E}\Delta_x^{k+1},
% \qquad q_k:=\mathbb{E}\|e^k\|^2.
% \]
% Adding this inequality with weight $a_0=(8\eta_1L_F^2)^{-1}$ to
% $\mathbb{E}(\Psi_{k+1}-\Psi_k)\leq
% -(2\eta_1)^{-1}\mathbb{E}\Delta_x^{k+1}+\eta_kq_k/2$ gives
% \[
% (\mathbb{E}\Psi_{k+1}+a_0q_{k+1})
% -(\mathbb{E}\Psi_k+a_0q_k)
% \leq-\frac{\rho_k}{16\eta_1L_F^2}q_k
% +\frac{\rho_k^2\sigma^2}{4\eta_1L_F^2b_{k+1}},
% \]
% when $\rho_k\geq8\eta_1\eta_kL_F^2$.
% Thus no division by $1-\rho_k$ or sign assumption on a difference of errors
% is needed. For the actual two-block oracle, however, (C.1) also contains
% $\Delta_B^{k+1}$, which is handled next.

% \smallskip
% \noindent\textbf{C.4. A complete absorption argument for both blocks.}
% Fix a horizon $T\geq1$ and a batch size $b\geq1$. All parameters below
% are constant over the iterations of this run. Set $\delta=1/6$ and
% \[
% \eta_T=\frac{1}{16L_FT^{1/3}},\qquad
% \tau_T=\frac{\eta_T}{2},\qquad
% \rho_T=\frac{16\eta_T^2L_F^2}{b},\qquad
% b_1=\lceil c_bT^{1/3}\rceil,\quad c_b>0.
% \tag{C.2}
% \]
% Abbreviate these parameters by $\eta,\tau,\rho$. Then $0<\rho\leq1/16$,
% $\tau L_F\leq1/32$, and
% \[
% \begin{aligned}
% \widetilde L_F&=L_F+\tfrac{13}{2}\tau L_F^2,\qquad
% \eta\widetilde L_F\leq\tfrac{77}{1024}<\tfrac12,\\
% C_x&=\eta^{-1}-\widetilde L_F\geq\frac1{2\eta},\\
% C_B&=\frac1{2\tau}-L_F-\frac{13}{4}\tau L_F^2
% \geq\frac1{4\tau}=\frac1{2\eta}.
% \end{aligned}
% \]
% For the last inequality, put $t=\tau L_F\leq1/32$ and use
% $t+(13/4)t^2\leq141/4096<1/4$.
% The correct pathwise StoRM Lyapunov descent is
% \[
% \Psi_{k+1}-\Psi_k
% \leq-C_B(\Delta_B^{k+1}+\Delta_B^k)
% -C_x\Delta_x^{k+1}+\frac\eta2\|e^k\|^2,
% \quad \Psi_k:=\Psi_r(\mathcal{Q}^k).
% \tag{C.3}
% \]
% The term $-(C_4-K_6)\Delta_x^k$ in the submitted (D.2) must be deleted:
% $\Psi_r$ has no $C_4\Delta_x^k$ term. Equation (C.3) follows directly
% from the preceding one-step bound and the definition of $\Psi_r$.

% Since $(1-\rho)^2\leq1-\rho$ and $(1-\rho)^2\leq1$, (C.1) implies
% \[
% q_{k+1}-q_k\leq-\rho q_k+\frac{2\rho^2\sigma^2}{b}
% +\frac{4L_F^2}{b}\mathbb{E}[\Delta_x^{k+1}+\Delta_B^{k+1}].
% \]
% Define $V_k=\mathbb{E}\Psi_k+(\eta/\rho)q_k$. Add the last inequality
% with weight $\eta/\rho$ to the full expectation of (C.3). Using
% $4\eta L_F^2/(b\rho)=1/(4\eta)$ gives
% \[
% \begin{aligned}
% V_{k+1}-V_k\leq{}&-\frac{C_B}{2}\mathbb{E}\Delta_B^{k+1}
% -C_B\mathbb{E}\Delta_B^k
% -\frac1{4\eta}\mathbb{E}\Delta_x^{k+1}\\
% &-\frac\eta2q_k+\frac{2\eta\rho\sigma^2}{b}.
% \end{aligned}
% \tag{C.4}
% \]
% This argument uses only nonnegative weights and does not require
% $\mathbb{E}\Psi_k$ to be nonincreasing.

% \smallskip
% \noindent\textbf{C.5. Stationarity and oracle complexity.}
% The blockwise stationarity calculation gives
% \[
% \begin{aligned}
% r_k&:=\mathbb{E}\operatorname{dist}^2
% (0,\partial\mathcal{L}(x^k,B^k,\Lambda^k))\\
% &\leq S_1\mathbb{E}\Delta_B^{k+1}
% +S_2\mathbb{E}\Delta_B^k
% +S_3\mathbb{E}\Delta_x^{k+1}+2q_k,
% \end{aligned}
% \]
% with
% \[
% S_1=4\tau^{-2}+12,\quad
% S_2=4[1+3(1+\tau L_F)^2],\quad
% S_3=4[\eta^{-2}+L_F^2+3(\tau L_F)^2].
% \]
% Use the StoRM value of $S_3$ derived in (D.9), rather than the StoM
% coefficient cited in that lemma's statement. Also, with $e^k=D^k-g(z^k)$,
% the identity for the $x$-component uses
% $g(z^k)=\eta^{-1}(x^k-x^{k+1})-e^k$; its squared-norm bound is unchanged.
% Put
% \[
% M_T=\max\left\{\frac{2S_1}{C_B},\frac{S_2}{C_B},
% 4\eta S_3,\frac4\eta\right\}.
% \]
% By (C.4),
% $r_k\leq M_T[V_k-V_{k+1}+2\eta\rho\sigma^2/b]$.
% Assume a finite lower bound $\mathbb{E}\Psi_k\geq\Psi_{*,T}$ and a
% uniform expected initial gap
% \[
% 0\leq\Delta_T:=\mathbb{E}\Psi_1-\Psi_{*,T}\leq\overline\Delta<\infty
% \quad\text{for all }T.
% \tag{C.5}
% \]
% Together with $q_1\leq\sigma^2/b_1$, telescoping yields
% \[
% \frac1T\sum_{k=1}^{T}r_k
% \leq\frac{M_T}{T}
% \left(\Delta_T+\frac{\eta\sigma^2}{\rho b_1}\right)
% +\frac{2M_T\eta\rho\sigma^2}{b}.
% \tag{C.6}
% \]
% Here $M_T=O(T^{1/3})$: the four entries of the maximum are respectively
% $O(T^{1/3})$, $O(T^{-1/3})$, $O(T^{1/3})$, and $O(T^{1/3})$.
% Substituting (C.2) and (C.5) into (C.6) gives
% \[
% \frac1T\sum_{k=1}^{T}r_k=O(T^{-2/3}).
% \]
% Uniform output selection over exactly these indices therefore attains
% expected squared stationarity error at most $\varepsilon^2$ with
% $T=O(\varepsilon^{-3})$. The fresh initial batch and two gradient
% evaluations per subsequent sample cost at most $b_1+2bT=O(T)$ SFO calls.
% The constants may depend on $L_F,b,c_b,\overline\Delta$, but not on $T$.

% For comparison with the original proof, the corrected coefficient in
% (D.15) is
% \[
% \frac{\widetilde S_4}{C_B}
% +\frac{2(C_\delta+4)}{\eta_1},
% \qquad \widetilde S_4=\max\{S_1,S_2\},
% \]
% without the extra $\widetilde S_4$ in the second denominator. In obtaining
% this bound, enlarge coefficients only on the nonnegative quantity
% $\Psi_k-\Psi_{k+1}+\eta\|e^k\|^2/2$, before separating its terms.
% Equation (C.6) supplies a direct replacement for the subsequent argument.

% \smallskip
% \noindent\textbf{C.6. Scope and prerequisites for the revised manuscript.}
% The batch-index correction in (4.8) is confirmed. The schedule (C.2) is
% one sufficient way to retain the complexity order after accounting for
% both blocks; it changes the submitted stipulation that $\tau$ is
% independent of $T$. No claim is made that every rate-preserving repair
% must use this schedule. The finite moments and uniform expected gap in
% (C.5) must be supplied by the boundedness assumptions or their proof;
% almost-sure boundedness alone does not establish them.

% A sufficient explicit alternative is to assume that all iterates and the
% warmup states lie almost surely in one deterministic compact set,
% uniformly over $T$. On this set, let $\|x^k\|\leq R$ and
% $\|\nabla_BF(x^{k+1},B^k)+\Lambda^k\|\leq M_B$.
% The $B$-update gives $\Delta_B^{k+1}\leq\tau^2M_B^2/2$.
% For the W-regularizer, $h^*$ is finite and continuous on its compact box
% domain, so $\mathcal{L}$ has finite bounds $L_{\min},L_{\max}$ on this set.
% Since $C_1=1/(2\tau)$, $C_3=\tau L_F^2/2$, and
% $C_2=1/(2\tau)-L_F+11\tau L_F^2/4$, one obtains
% \[
% \Psi_k\geq L_{\min}-\frac{\tau M_B^2}{4}-\tau L_F^2R^2,
% \qquad
% \Psi_1\leq L_{\max}+\frac{\tau M_B^2}{4}
% +\frac{11\tau^3L_F^2M_B^2}{8}.
% \]
% Thus (C.5) follows under this stronger, explicitly uniform boundedness
% assumption. It must not be conflated with sample-path boundedness with
% an uncontrolled random radius.
% The first analyzed index must have a valid preceding $B$- and dual
% update. One explicit initialization is a single $B$- and dual warmup
% at fixed $x$, followed by a fresh batch for $D^1$; then reindex the first
% analyzed state as $1$. This ensures both the dual subgradient identity
% and the initial variance bound. The output range and summations must
% use the same reindexing. These are conditions of the repaired result,
% not conclusions inferred from the submitted initialization.

% \clearpage
% \section*{Appendix D: A Sufficient Repair with Fixed $\tau$}
% \noindent\emph{Internal working notes. This argument retains a positive
% $\tau$ independent of $T$ and the corrected upper bound from Lemma~C.2.
% It adds an explicit block-stationary initialization and a tracking estimate.
% The horizon-dependent choice in Appendix~C remains an alternative.}

% \smallskip
% \noindent\textbf{D.1. Model structure and initialization.}
% Use the quadratic coupling in Problem~(P), and put
% \[
% m:=\frac\gamma n>0,\qquad U^k:=\mathcal{D}(x^k),\qquad
% \nabla_BF(x,B)=m(B-\mathcal{D}(x)).
% \]
% The smoothness bound satisfies $m\leq L_F$. Fix $\delta=1/6$ and
% \[
% 0<\tau<\frac{\sqrt{30}-2}{13L_F},\qquad
% r:=\tau m\leq\tau L_F<\frac27.
% \tag{D.1}
% \]
% At the given initial network parameter $x^1$, choose
% \[
% B^1=\operatorname{prox}_{m^{-1}h^{**}}(U^1),\qquad
% \Lambda^1=m(U^1-B^1),\qquad B^0=B^1,\quad\Lambda^0=\Lambda^1.
% \tag{D.2}
% \]
% Then $m(B^1-U^1)+\Lambda^1=0$ and
% $B^1\in\partial h^*(\Lambda^1)$. Thus both preceding update identities
% used in (5.9) hold at the initial index. Initialize $D^1$ from a fresh
% batch of size $b_1$.
% For the W-regularizer, $h^{**}(B)=\lambda\sum_{ij}(|B_{ij}|-1)_+$,
% so (D.2) is explicit, componentwise:
% \[
% [\operatorname{prox}_{m^{-1}h^{**}}(U)]_{ij}
% =\begin{cases}
% U_{ij},& |U_{ij}|\leq1,\\
% \operatorname{sgn}(U_{ij}),&1<|U_{ij}|\leq1+\lambda/m,\\
% U_{ij}-(\lambda/m)\operatorname{sgn}(U_{ij}),&|U_{ij}|>1+\lambda/m.
% \end{cases}
% \]
% In particular, if $|U_{ij}^1|\leq1$ for every entry, the choice reduces to
% $B^1=U^1$ and $\Lambda^1=0$. This output restriction is not assumed in
% the general argument.

% \smallskip
% \noindent\textbf{D.2. Increment dynamics for the deterministic block.}
% The dual step here is the minimization form of (4.5); its equivalent
% proximal formula must include the factor $\tau^{-1}$ before
% $2B^{k+1}-B^k$. With $Z^k=\tau\Lambda^k$, define the scaled proximal map
% \[
% P(v)=\tau\operatorname{prox}_{\tau^{-1}h^*}(v/\tau).
% \]
% It is coordinatewise, monotone and $1$-Lipschitz. The updates read
% \[
% B^{k+1}=(1-r)B^k-Z^k+rU^{k+1},\qquad
% Z^{k+1}=P(Z^k+2B^{k+1}-B^k).
% \]
% Define
% \[
% \beta_k=B^k-B^{k-1},\quad
% \zeta_k=Z^k-Z^{k-1},\quad
% w_k=\zeta_k-\beta_k,\quad
% u_k=U^{k+1}-U^k.
% \]
% For each coordinate, the secant slope of $P$ between the two consecutive
% arguments is some $d\in[0,1]$. Subtracting the updates yields
% \[
% \begin{pmatrix}\beta_{k+1}\\w_{k+1}\end{pmatrix}
% =A_d\begin{pmatrix}\beta_k\\w_k\end{pmatrix}+f_du_k,
% \quad
% A_d=\begin{pmatrix}-r&-1\\r(1-2d)&1-d\end{pmatrix},\quad
% f_d=\begin{pmatrix}r\\r(2d-1)\end{pmatrix}.
% \tag{D.3}
% \]
% The slope may differ across coordinates and iterations. No differentiability
% of the proximal map is required.

% \smallskip
% \noindent\textbf{D.3. A common quadratic estimate.}
% Let
% \[
% \Pi_r=\begin{pmatrix}r&r\\r&1\end{pmatrix}\succ0,\qquad
% H_k=r\|\beta_k\|_F^2+2r\langle\beta_k,w_k\rangle+\|w_k\|_F^2.
% \]
% For $0<r\leq2/7$, direct calculation gives
% \[
% \begin{aligned}
% \Pi_r-A_0^\top \Pi_rA_0
% &=\begin{pmatrix}r-r^2+r^3&r^2\\r^2&r\end{pmatrix}
% \succeq\frac r2 I,\\
% \Pi_r-A_1^\top \Pi_rA_1
% &=\begin{pmatrix}r-r^2-3r^3&r-2r^2\\r-2r^2&1-r\end{pmatrix}
% \succeq\frac r4 I.
% \end{aligned}
% \]
% For the first inequality, the shifted leading entry is positive and its
% determinant is $r^2(1/4-r/2-r^2/2)\geq13r^2/196$.
% For the second inequality, the shifted leading entry divided by $r$ is
% $3/4-r-3r^2\geq43/196$, and the shifted determinant is $rp(r)$, where
% \[
% p(r)=\frac34-\frac{47}{16}r+\frac94r^2-\frac14r^3
% \geq p(2/7)=\frac{243}{2744}>0.
% \]
% The function $p$ decreases on this interval.
% Since $A_d=(1-d)A_0+dA_1$, convexity gives, for every vector $y$,
% \[
% \|A_dy\|_{\Pi_r}^2
% \leq(1-d)\|A_0y\|_{\Pi_r}^2+d\|A_1y\|_{\Pi_r}^2
% \leq\|y\|_{\Pi_r}^2-\frac r4\|y\|^2.
% \]
% Thus $A_d^\top\Pi_rA_d\preceq\Pi_r-rI/4$ for every $d\in[0,1]$.
% Moreover,
% $\|\Pi_r\|\leq9/7$, $\|A_d\|\leq3/2$, and
% $\|f_d\|\leq\sqrt2r$, so
% $\|A_d^\top \Pi_rf_d\|\leq3r$ and
% $f_d^\top \Pi_rf_d\leq r$.
% Using $6r\|y\||u|\leq(r/8)\|y\|^2+72r|u|^2$ in (D.3), and summing
% across coordinates, yields
% \[
% H_{k+1}-H_k
% \leq-\frac r8(\|\beta_k\|_F^2+\|w_k\|_F^2)
% +73r\|U^{k+1}-U^k\|_F^2.
% \tag{D.4}
% \]
% Also,
% $\|\beta_{k+1}\|_F^2\leq
% 3(\|\beta_k\|_F^2+\|w_k\|_F^2)+3r^2\|u_k\|_F^2$.
% Combining with (D.4) and the Lipschitz property of $\mathcal{D}$ gives
% \[
% \Delta_B^{k+1}
% \leq\frac{12}{r}(H_k-H_{k+1})+G\Delta_x^{k+1},
% \qquad G:=1753L_{\mathcal{D}}^2.
% \tag{D.5}
% \]
% The numerical constants are sufficient, not optimized. The difference
% $H_k-H_{k+1}$ need not be nonnegative; (D.5) follows by algebraic
% rearrangement, with no sign assumption on that difference.

% \smallskip
% \noindent\textbf{D.4. Combining the stochastic and tracking estimates.}
% Take constant $\eta_k=\eta_T$ and $\rho_k=\rho_T$ within each run. Write
% $q_k=\mathbb{E}\|e^k\|^2$, $\Psi_k=\Psi_r(\mathcal{Q}^k)$, and
% $d_0=48L_F^2/(br)$. The corrected same-batch recursion (C.1) and (D.5)
% imply
% \[
% \begin{aligned}
% &(q_{k+1}+d_0\mathbb EH_{k+1})-(q_k+d_0\mathbb EH_k)\\
% &\qquad\leq-\rho_Tq_k+\frac{2\rho_T^2\sigma^2}{b}
% +\frac{4L_F^2(1+G)}{b}\mathbb{E}\Delta_x^{k+1}.
% \end{aligned}
% \]
% Retain the original powers in the parameter schedule:
% \[
% \eta_T=\frac{\eta_0}{\widetilde L_FT^{1/3}},\qquad
% \rho_T=\frac{8\kappa\eta_0^2}{T^{2/3}},\qquad
% b_1=\lceil c_bT^{1/3}\rceil,
% \tag{D.6}
% \]
% where $\widetilde L_F=L_F+(13/2)\tau L_F^2$ and the numerical constants satisfy
% \[
% \kappa\geq\max\left\{1,\frac{2L_F^2(1+G)}{b\widetilde L_F^2}\right\},
% \qquad 0<\eta_0\leq\min\left\{\frac12,\frac1{4\sqrt\kappa}\right\}.
% \tag{D.7}
% \]
% Thus $\eta_T\widetilde L_F\leq1/2$, $0<\rho_T\leq1/2$, and
% $4\eta_TL_F^2(1+G)/(b\rho_T)\leq1/(4\eta_T)$.
% With
% \[
% W_k=\mathbb{E}\Psi_k+\frac{\eta_T}{\rho_T}
% (q_k+d_0\mathbb EH_k),
% \]
% adding the last error estimate to (C.3) gives
% \[
% \begin{aligned}
% W_{k+1}-W_k\leq{}&-C_B\mathbb{E}[\Delta_B^{k+1}+\Delta_B^k]
% -\frac1{4\eta_T}\mathbb{E}\Delta_x^{k+1}\\
% &-\frac{\eta_T}{2}q_k+\frac{2\eta_T\rho_T\sigma^2}{b}.
% \end{aligned}
% \tag{D.8}
% \]
% This avoids the invalid coefficient replacement in the original (D.7)
% and accounts explicitly for the omitted $B$-increment.

% \smallskip
% \noindent\textbf{D.5. Rate, initialization, and scope.}
% Let $S_1,S_2,S_3$ be the StoRM stationarity coefficients in Appendix~C.5,
% and put
% \[
% N_T=\max\left\{\frac{S_1}{C_B},\frac{S_2}{C_B},
% 4\eta_TS_3,\frac4{\eta_T}\right\}=O(T^{1/3}).
% \]
% The block-stationary initialization (D.2) gives $H_1=0$. Under finite
% moments, choose a finite lower bound
% $\mathbb{E}\Psi_k\geq\Psi_{*,T}$ and assume the uniform initial gap
% $0\leq\Delta_T:=\mathbb{E}\Psi_1-\Psi_{*,T}\leq\overline\Delta$.
% Telescoping (D.8) gives
% \[
% \frac1T\sum_{k=1}^{T}\mathbb{E}\operatorname{dist}^2
% (0,\partial\mathcal{L}(x^k,B^k,\Lambda^k))
% \leq\frac{N_T}{T}
% \left(\Delta_T+\frac{\eta_T\sigma^2}{\rho_Tb_1}\right)
% +\frac{2N_T\eta_T\rho_T\sigma^2}{b}
% =O(T^{-2/3}).
% \tag{D.9}
% \]
% The SFO complexity order is $O(\varepsilon^{-3})$, with aligned averaging
% and output indices. Constants may depend on $\tau,m, L_{\mathcal{D}}$ and
% other fixed problem parameters, including $n$, but not on $T$.
% Computing $U^1=\mathcal{D}(x^1)$ and the entrywise initialization incurs a
% one-time deterministic cost, which should be stated separately from
% stochastic gradient calls.

% For a general initial tracking energy, the same proof retains the extra
% term
% \[
% \frac{N_T\eta_Td_0\mathbb EH_1}{\rho_TT}=O(T^{-1/3})\mathbb EH_1.
% \]
% This term limits this particular upper bound. It is not a lower bound
% on the algorithm's actual error and does not prove that fixed $\tau$
% fails under every other initialization or every other proof approach.
% The sufficient result above adds the explicit initialization (D.2) and
% numeric restrictions (D.7), and uses the Lipschitz mapping assumption
% already discussed in the manuscript's Appendix~B. These requirements
% must be stated in the revised theorem. The expectation-gap condition
% remains the separate boundedness requirement.

\end{document}
